\begin{theorem}[G-25: Exact strict-direction comparison through $2m-2$ rows]\label{G-25}
Every assertion and every explicit bound of G-23 remains valid when its
row range $4\le r\le m$ is replaced by
\[
             4\le r\le2m-2,
\]
with $n>m\ge17$ and the same strict conditions
$\nu_r>0$, $\nu_1>(n-m)s$, $\nu\vdash ns$.
In particular the exact eventual exponent is still
\[
 \Delta=rm^2-(r-2)n^2-2m+r-1\ne0,
\]
and both multiplicities are positive at every integer stretch
$j\ge\beta^2$, with the same definitions of $b,L,p,\beta,T,E,u,v$.
The effective thresholds for strict comparison, and the exact
floor boundary $N_{m,r}$, are those of G-23.

In the newly added range $m<r\le2m-2$, the exponent is positive
if and only if
\[
                   r=n=m+1.
\]
For these parameters it equals one. Thus, for every $m\ge17$,
the explicit nonrectangular direction
\[
 s=2,\qquad \nu=(3,2^{m-1},1),\qquad n=r=m+1
\]
has an eventually positive genuine-padding multiplicity gap in
every sufficiently large coordinate degree $2j$, and its
permanent-to-determinant multiplicity ratio is
$\kappa_\nu j(1+O(j^{-1}))$, with $\kappa_\nu>0$.
Every other pair of integer parameters in this newly added row
range has eventual determinant dominance in every fixed strict
direction.
\end{theorem}

\begin{proof}
R-29 gives both the common stable tensor and a common stable tensor
with no nontrivial scalar multiple in its orbit throughout the
enlarged range. Its enhanced-coordinate extraction agrees exactly
with that in G-22. The dimension formulas used by G-23 already
hold here: G-21 proves the permanent formula for every $r\ge4$
when $m\ge17$, and R-28 proves the determinant formula and the
padding dimension for every $r\ge4$ in the present size range.
Consequently the invariant-domain dimensions are precisely the
same $u,v$ displayed in G-23.

The only numerical estimate in that proof that used $r\le m$
beyond the stable-point input was $u,v\ge2$. We supply the
replacement explicitly. Since $r-2\ge r/2$, $r\le2m-2$,
$n\ge m+1$ and $r_\eta\ge0$,
\[
 \begin{aligned}
 u&=(r-2)n^2-r^2+r_\eta+3
       \ge\frac r2(m^2-2m+5)+3>2,\\
 v&=r(m^2-r+1)-2m+r_\eta+2
       \ge r(m^2-2m+3)-2m+2>2.
 \end{aligned}
\]
The last inequality follows already from $r\ge4$ and $m\ge17$.

All subsequent steps of G-23 use only these exact dimensions,
the ambient representation in rank $r$, and the two properties
of the common tensor supplied by R-29. Derksen's action-degree
and generator bounds are unchanged. The same nonzero generator
degrees have gcd one by closed-orbit separation, so the common
invariants $h_j$ exist for every $j\ge\beta^2$. The same
generator grouping gives the standard $E$th Veroneses, and the
same degree-$p$ parameter maps give their image-degree bounds.
Their every-residue Hilbert sandwiches, the Noether lower bound,
and the Chardin upper bound therefore prove the stated effective
thresholds and ratio with these actual values of $u,v$.
The parity and nonsquare argument excluding $\Delta=0$ in G-23
was valid for every integer $r\ge4$, with no upper row bound.
This proves the full extension, including all its constants.

For the additional sign classification, fix $m<r$.
The exponent is strictly decreasing as an integer function of
$n\ge m+1$, and direct substitution gives
\[
       \Delta\big|_{n=m+1}=2m(m-r+1)+1.
\]
If $r\ge m+2$, this is at most $1-2m<0$.
If $r=m+1$, its value at $n=m+1$ is one, while increasing $n$
to $m+2$ subtracts $(m-1)(2m+3)>1$. This proves the asserted
unique positive parameter case. Finally the displayed partition
has $m+1$ positive parts and total size $2m+2=ns$; its first
part is $3>2=(n-m)s$. It therefore meets the strict hypotheses,
and substituting $\Delta=1$ proves its claimed ratio and gap.
\end{proof}

\paragraph{Scope.}
The new positive example specifies a total multiplicity inequality,
not merely one nonzero equation copy. It is an eventual statement
with the explicit G-23 threshold and does not assert the comparison
at every earlier degree. No rectangle satisfies the strict
condition when $r>m$ and $n>m$ are integers: then
$rm/(r-1)\le m+1$. The rectangular boundary result G-24 keeps
its stated range $r\le m$. General nonrectangular equality
faces, more than $2m-2$ rows, changing directions, and
polynomial-resource coefficient extraction remain separate gates.
